% ===== PRÉAMBULE CANONIQUE — à coller VERBATIM en tête de CHAQUE .tex =====
\documentclass[11pt,a4paper]{article}
\usepackage[utf8]{inputenc}\usepackage[T1]{fontenc}\usepackage{lmodern}
\usepackage[french]{babel}
\usepackage{amsmath,amssymb,mathtools}\usepackage{icomma}
\usepackage[version=4]{mhchem}          % \ce{...} : équations & formules
\usepackage{siunitx}\sisetup{locale=FR} % \qty{}{}, \si{}, \ang{}
\usepackage{chemfig}                    % \chemfig{...} : structures développées
\usepackage{xcolor}\usepackage{colortbl}
\usepackage{tikz}\usepackage{pgfplots}\pgfplotsset{compat=1.18}
\usepackage[most]{tcolorbox}\usepackage{enumitem}\usepackage{array,booktabs}
\usepackage[a4paper,margin=2.2cm]{geometry}
\usetikzlibrary{arrows.meta,calc,angles,quotes,positioning,patterns,backgrounds,shapes.geometric}
\definecolor{primary}{RGB}{222,49,99}\definecolor{primarydark}{RGB}{150,22,55}
\definecolor{defcol}{RGB}{0,114,178}\definecolor{methcol}{RGB}{52,92,148}
\definecolor{excol}{RGB}{0,135,90}\definecolor{attncol}{RGB}{200,40,55}
\tcbset{boxbase/.style={enhanced,breakable,boxrule=0.9pt,arc=2.5pt,
  left=3mm,right=3mm,top=2.3mm,bottom=2.3mm,fonttitle=\bfseries,coltitle=white}}
% --- boîtes TUTORIEL (noms EXACTS, ne pas renommer) ---
\newtcolorbox{cle}[1][]{boxbase,colback=defcol!6,colframe=defcol,
  title={\textsc{À retenir}\ifx\relax#1\relax\relax\else\textmd{ : \itshape #1}\fi}}
\newtcolorbox{methode}[1][]{boxbase,colback=methcol!7,colframe=methcol,sharp corners,
  title={\textsc{Méthode}\ifx\relax#1\relax\relax\else\textmd{ --- \itshape #1}\fi}}
\newtcolorbox{exemplebox}[1][]{boxbase,colback=excol!5,colframe=excol,
  title={\textsc{Exemple}\ifx\relax#1\relax\relax\else\textmd{ : \itshape #1}\fi}}
\newtcolorbox{piege}[1][]{boxbase,colback=attncol!6,colframe=attncol,
  title={\textsc{Piège}\ifx\relax#1\relax\relax\else\textmd{ : \itshape #1}\fi}}
\newtcolorbox{remarque}[1][]{boxbase,colback=black!4,colframe=black!45,
  title={\textsc{Remarque}\ifx\relax#1\relax\relax\else\textmd{ : \itshape #1}\fi}}
% --- boîtes EXERCICE / CORRIGÉ (noms EXACTS, auto-comptées) ---
\newtcolorbox[auto counter]{exo}[1][]{boxbase,colback=primary!4,colframe=primary,
  title={\textsc{Exercice}~\thetcbcounter\ifx\relax#1\relax\relax\else\textmd{ --- \itshape #1}\fi}}
\newtcolorbox[auto counter]{corrige}[1][]{boxbase,colback=excol!4,colframe=excol,
  title={\textsc{Corrigé}~\thetcbcounter\ifx\relax#1\relax\relax\else\textmd{ --- \itshape #1}\fi}}
\newcommand{\R}{\mathbb{R}}\newcommand{\N}{\mathbb{N}}\newcommand{\Z}{\mathbb{Z}}\newcommand{\ds}{\displaystyle}
\begin{document}
\begin{exo}[Nommer les alcanes linéaires]
Donne le nom de l'alcane linéaire possédant :
\begin{enumerate}[label=\alph*)]
  \item $1$ carbone
  \item $4$ carbones
  \item $6$ carbones
  \item $8$ carbones
\end{enumerate}
\end{exo}

\begin{exo}[De la semi-développée à la formule brute]
Détermine la formule brute (compte tous les atomes) de chaque molécule.
\begin{enumerate}[label=\alph*)]
  \item \ce{CH3-CH2-CH2-CH3}
  \item \ce{CH3-CH2-OH}
  \item \ce{CH3-CO-CH3}
\end{enumerate}
\end{exo}

\begin{exo}[Combien d'hydrogènes sur ce carbone ?]
Dans un composé saturé, un carbone est lié à un certain nombre d'\emph{autres carbones}. Combien d'atomes d'hydrogène porte-t-il alors ? (Rappel : $n_{\text{H}}=4-\text{liaisons non-H}$.)
\begin{enumerate}[label=\alph*)]
  \item lié à $1$ autre carbone
  \item lié à $2$ autres carbones
  \item lié à $3$ autres carbones
  \item lié à $4$ autres carbones
\end{enumerate}
\end{exo}

\begin{exo}[Nommer la classe d'un carbone]
Selon le nombre d'autres carbones auxquels il est lié, donne la classe (primaire, secondaire, tertiaire ou quaternaire) du carbone.
\begin{enumerate}[label=\alph*)]
  \item lié à $1$ autre carbone
  \item lié à $2$ autres carbones
  \item lié à $3$ autres carbones
  \item lié à $4$ autres carbones
\end{enumerate}
\end{exo}

\begin{exo}[Classer les carbones du propane]
On considère le propane \ce{CH3-CH2-CH3}. Classe chacun de ses trois carbones (de gauche à droite) en primaire, secondaire, tertiaire ou quaternaire.
\end{exo}

\end{document}
